%% =====================================================================
%% IEEE GRSL submission source — paper 4
%% Compile: pdflatex -> bibtex (or manual thebibliography) -> pdflatex x2
%% Figures: figures/grsl/*.pdf  (place alongside or adjust \graphicspath)
%% =====================================================================
\documentclass[journal]{IEEEtran}

\usepackage{amsmath}
\usepackage{amssymb}
\usepackage{graphicx}
\usepackage{booktabs}
\usepackage{multirow}
\usepackage{url}
\usepackage[hidelinks]{hyperref}

%% Two search paths so the same file compiles both in the working tree
%% (../figures/grsl) and after unzipping the submission package (figures/grsl).
\graphicspath{{figures/grsl/}{../figures/grsl/}}

%% The two standard compatibility shims
\interdisplaylinepenalty=2500

%% Tighten float placement so the paper fits five pages without cutting content.
%% This is lossless: it only reclaims whitespace around figures.
\setlength{\textfloatsep}{7pt plus 2pt minus 2pt}
\setlength{\floatsep}{7pt plus 2pt minus 2pt}
\setlength{\intextsep}{7pt plus 2pt minus 2pt}
\setlength{\dbltextfloatsep}{7pt plus 2pt minus 2pt}
\setlength{\dblfloatsep}{7pt plus 2pt minus 2pt}
\renewcommand{\topfraction}{0.95}
\renewcommand{\bottomfraction}{0.95}
\renewcommand{\textfraction}{0.05}
\renewcommand{\floatpagefraction}{0.92}
\renewcommand{\dbltopfraction}{0.95}
\renewcommand{\dblfloatpagefraction}{0.92}

\begin{document}

\title{Test-Set Geographic Composition Dominates Method Ranking in
Remote Sensing Super-Resolution Evaluation}

%% Author block. Dawei Liu (advisor) is the corresponding author and carries the
%% asterisk. All three authors share one affiliation. The affiliation and the
%% corresponding-author contact are printed at the end of the manuscript (last
%% page) rather than as a title-page footnote, following the authors' request.
\author{Zhengyang~Chen,~Dawei~Liu\IEEEauthorrefmark{1},~and~Zhiheng~Liu}

\maketitle

\begin{abstract}
Remote sensing super-resolution is evaluated on a small number of geographically
distinct sites, and methods are ranked by differences of one or two tenths of a
decibel. We ask how much of such a ranking is decided by which sites happen to be
in the test set. Three multi-temporal Sentinel-2 networks, three seeds each, are
evaluated on all 18 sites of the S2DS benchmark separately, and every possible
$k$-site test subset is enumerated. The two strongest networks differ by 0.16~dB
on average and are statistically indistinguishable at the site level (Wilcoxon
$p = 0.32$, bootstrap 95\% CI $[-0.26, +0.59]$~dB). Over the 8568 five-site
subsets the better network wins in 67.6\% of them, and the benchmark's own three
partitions produce three different orderings of the same three architectures.
Site variance exceeds model variance by 31.7 times, and the PSNR range across
sites is 11.38~dB, which is 57 times the gap being resolved. A stable ranking
requires at least 16 sites. The effect replicates on MuS2, a public benchmark
with a different sensor pair and scale factor, where zero-training baselines
still span 6.31~dB across seven tiles.
\end{abstract}

\begin{IEEEkeywords}
Super-resolution, evaluation protocol, multi-temporal Sentinel-2, benchmark
design, geospatial sampling.
\end{IEEEkeywords}

%% ---------------------------------------------------------------- I
\section{Introduction}

Super-resolution (SR) for remote sensing is evaluated the way natural-image SR
is evaluated \cite{deepsum2020,diffusionsat2024}. A fixed test split is assembled,
every method is run over it, and a single mean PSNR is reported. The architecture
that wins by a tenth of a decibel is described as the better one.

This convention carries an assumption that is rarely examined, namely that the
test split is large and varied enough for its particular composition not to decide
the outcome. In natural-image SR the assumption is defensible, because benchmark
splits contain hundreds or thousands of images drawn from the same photographic
distribution. Remote sensing data are different. A benchmark for satellite SR is
assembled from a small number of geographical sites \cite{mus2}. Each site is
internally homogeneous, and the sites differ from one another in climate, terrain,
land cover and atmospheric conditions. The number of independent geographical
units is often five or fewer.

Two properties follow. Images from one site are spatially correlated, so samples
drawn from a single site are not independent draws from a distribution over
scenes. The difficulty of the task itself varies with the site, because texture,
contrast and cloud frequency are governed by geography. Both properties push in
the same direction, and the number of sites determines how much of that dependence
is averaged out.

We test this directly. Rather than treating the benchmark as a single number, we
evaluate every model on every site separately and ask what happens to the rankings
as the site composition of the test set changes. The experiment uses nine existing
checkpoints.

The contributions are the following. We quantify site-level dispersion in a
multi-temporal Sentinel-2 SR benchmark, where the PSNR across 18 sites spans
11.38~dB while the two strongest networks differ by 0.16~dB. We decompose variance
between sites and between models, finding a ratio of 31.7. We enumerate every
five-site test subset and show that the ranking of the two strongest networks
flips in 32.9\% of them, and that at least 16 sites are needed for stability. We
replicate the effect on MuS2, a benchmark built by others with a different sensor
pair and scale factor. Finally, we identify a second way a benchmark can mislead,
in that the fraction of valid target pixels falls as low as 1.5\% in one site.

%% ---------------------------------------------------------------- II
\section{Data and Methods}

\subsection{S2DS Benchmark}

S2DS is a multi-temporal SR benchmark built from real Sentinel-2 Level-2A
acquisitions \cite{drusch2012,s2psd}. It contains 18 sites partitioned by geography
into 10 training, 3 validation and 5 test sites. Each sample provides $T = 12$
low-resolution frames acquired on distinct dates at 40~m, the bands B02, B03, B04
and B08, per-frame per-pixel cloud fraction and per-frame signed time offsets,
together with a 10~m target image for the reference date. The scale factor is 4.
The split is by site and not by patch, so training and test samples come from
disjoint areas. The evaluated splits contain 873, 180 and 423 samples respectively.

\subsection{Models}

Three architecture families are used, each with three seeds. Mamba-MISR carries
inter-frame interaction in a selective state-space module \cite{mamba} and gates
fusion by per-pixel observation reliability, with temporal offsets entering the
gate only. HighRes-Net-cld is the recursive fusion architecture of
\cite{deudon2020} extended with cloud-mask conditioning. BreizhSR is a lightweight
multi-frame network \cite{okabayashi2024} used as a lower-capacity reference. All
twelve checkpoints are fixed at 30\,000 iterations, with seeds 2026, 2027 and 2028.

\subsection{Evaluation Protocol}

PSNR is computed over the valid mask, defined as pixels whose target-date cloud
label is clear \cite{cmix}, following the standard image-quality measures of
\cite{ssim2004}. Per-site values are unweighted means over the samples of that
site. The two aggregation conventions differ, and we state which is used at each
point, because on the test split the sample-weighted mean is 27.2285~dB for
Mamba-MISR seed 2026 while averaging the five site means gives 27.0993~dB, a
difference of 0.1293~dB arising from unequal sample counts. Samples with an empty
valid mask are excluded rather than counted as zero, leaving 1314 of 1476 samples
for the network comparison.

\subsection{Analysis Design}

Every model--seed combination is evaluated separately on each of the 18 sites,
giving a $9 \times 18$ matrix of site-level PSNR. Variance across site means and
variance across model--seed means are computed from that matrix, and the ratio is
reported with a bootstrap 95\% confidence interval obtained by resampling both
sites and model--seed combinations over 2000 replicates.

For a pair of networks the win rate of the first is computed by exact enumeration
of every subset of size $k$ drawn from the 18 sites, for $k$ from 3 to 18. With
$\binom{18}{5} = 8568$ the enumeration is exact and no sampling approximation is
used. The paired site-level difference is tested with a Wilcoxon signed-rank test
\cite{wilcoxon1945}, with the rank-biserial correlation as effect size and a
bootstrap confidence interval over the 18 sites. The unit of analysis is the site
and not the sample, because samples within one site are not independent. Means and
medians are reported together, since the distribution of site-level differences is
skewed \cite{lakens2017}.

\subsection{External Replication on MuS2}

MuS2 \cite{mus2} is a public multi-image SR benchmark built by others from
Sentinel-2 as the low-resolution source and higher-resolution multispectral
imagery as the reference, with a scale factor of 3. The local copy holds 91 scenes
over 7 MGRS tiles.

Two settings are used and reported separately. In Setting~A the reference image is
downsampled by 3 and reconstructed with nearest, bilinear or bicubic interpolation.
Ground truth and prediction then live on the same radiometric scale, so no
registration or band-matching uncertainty enters, and the quantity measured is how
hard the fabric of a given tile is to reconstruct. In Setting~B a Sentinel-2 band
is upsampled and compared against the reference band, after a quantile match that
maps the prediction onto the reference scale. Registration between the two products
is not independently verified, so the Setting~B numbers carry that uncertainty and
are quoted alongside Setting~A rather than on their own.

%% ---------------------------------------------------------------- III
\section{Results}

\subsection{Site-Level Dispersion}

Fig.~\ref{fig:aoi} shows the valid PSNR of Mamba-MISR averaged over nine
model--seed combinations for each of the 18 sites. The values span 11.38~dB, from
18.56~dB at \texttt{qingdao\_inland} to 29.94~dB at \texttt{sjz\_north\_agri}, with
a standard deviation of 2.81~dB. The three partition memberships behave as one
population. The 10 training sites average 27.33~dB, the 3 validation sites
23.75~dB and the 5 test sites 26.27~dB, and the highest and lowest sites are not
separated by partition. \texttt{dunhuang\_oasis}, a test site, ranks third, while
\texttt{harbin\_agri}, a training site, ranks second from the bottom. Considered
per architecture rather than averaged, the spans across the same sites are 12.45,
13.08 and 10.96~dB, so the dispersion is not an artefact of averaging.

\begin{figure}[!t]
\centering
\includegraphics[width=\columnwidth]{fig1_aoi_psnr.pdf}
\caption{Valid PSNR at each of the 18 S2DS sites, averaged over nine model--seed
combinations and sorted by score. Colour marks the training, validation and test
partitions. Site identifiers follow Table~I. The 11.38~dB span across sites exceeds
the 0.16~dB gap between the two strongest networks by a factor of 71.}
\label{fig:aoi}
\end{figure}

A partition effect does exist but is far smaller. Evaluating Mamba-MISR seed 2026
on the training partition gives 27.4640~dB and on the test partition 27.2285~dB, a
difference of 0.2355~dB, which is within the range that single-seed training noise
produces on this benchmark.

\subsection{Variance Attribution}

Variance across site means is 7.886~dB$^2$, a standard deviation of 2.81~dB.
Variance across model--seed means is 0.249~dB$^2$, a standard deviation of 0.50~dB.
The ratio is 31.7 with a bootstrap 95\% confidence interval of 6.7 to 307. The
interval is wide because 18 sites and 9 model--seed combinations give limited
resolution at the upper end, but its lower bound already places site-to-site
variation an order of magnitude above model-to-model variation. The ratio depends
on which models are compared, and Section~\ref{sec:repl} shows a public benchmark
on which the same computation over four zero-training baselines returns a ratio
below one. What is stable across datasets is the site-level dispersion, not the
ratio.

\subsection{Ranking Instability}

Comparing Mamba-MISR against HighRes-Net-cld, their site-level means over 18 sites
are 26.6870~dB and 26.8475~dB. The paired difference has a mean of 0.1605~dB and a
median of 0.4474~dB. A Wilcoxon signed-rank test gives $p = 0.3247$ with a
rank-biserial correlation of 0.275, and the bootstrap 95\% confidence interval of
the mean difference is $[-0.2593, +0.5856]$~dB, which contains zero. The two
networks are statistically indistinguishable at the site level.

Fig.~\ref{fig:flip} shows the consequence. Enumerating all 8568 five-site subsets,
HighRes-Net-cld scores higher in 67.6\% of them, so Mamba-MISR wins 32.9\%. A study
drawing a different five sites would have reported the opposite ordering in one
case out of three. The comparison is informative because it sits at a specific gap
size. Two further pairs on the same matrix differ by far more. Mamba-MISR exceeds
BreizhSR by 0.90~dB and wins in 99.43\% of five-site subsets, and HighRes-Net-cld
exceeds BreizhSR by 1.06~dB and wins in 99.03\%. A gap near 0.16~dB produces a win
rate near one third, while gaps near one decibel produce rates above 99\%, which
locates the transition region.

The absolute level is equally sensitive. Drawing five sites at random, the mean
PSNR of the nine model--seed combinations ranges from 23.40 to 29.94~dB, with a
standard deviation of 1.14~dB and a span of 6.54~dB. Raising the subset size to ten
narrows the standard deviation to 0.62~dB.

A third view comes from the partition the benchmark defines, which was fixed before
this analysis. Ranking the three architectures within each partition gives three
different orderings, as Table~\ref{tab:main} shows. On the 10 training
sites HighRes-Net-cld leads with 28.20~dB against 27.29~dB. On the 5 test sites the
order reverses, with Mamba-MISR at 27.06~dB against 26.11~dB. On the 3 validation
sites Mamba-MISR leads at 24.07~dB while HighRes-Net-cld falls to last place at
23.56~dB, below BreizhSR at 23.64~dB. The validation and test partitions are
equally unseen by all three models, so the reversal between them is attributable
to which sites they contain and to nothing else.

\begin{figure}[!t]
\centering
\includegraphics[width=\columnwidth]{fig2_flip_curve.pdf}
\caption{Win rate of the first method as the number of test sites $k$ grows, by
exact enumeration of all $\binom{18}{k}$ subsets. The shaded band marks the region
in which the ranking is unstable. A gap of 0.16~dB leaves the ordering undecided
until $k = 16$, while gaps near one decibel are stable from $k = 3$.}
\label{fig:flip}
\end{figure}

\begin{table}[!t]
\caption{Valid PSNR (dB) at Each S2DS Site, Mean Over Three Seeds}
\label{tab:main}
\centering
\scriptsize
\setlength{\tabcolsep}{2.6pt}
\renewcommand{\arraystretch}{0.90}
\begin{tabular}{lcccc}
\toprule
Site & Split & MISR & HN-cld & BreizhSR \\
\midrule
sjz north agri   & Tr & 29.83 & 31.05 & 28.93 \\
xian plain       & Tr & 29.27 & 30.48 & 28.37 \\
chengdu plain    & Tr & 31.09 & 31.30 & 25.64 \\
lanzhou valley   & Tr & 29.16 & 29.85 & 28.94 \\
dunhuang oasis   & Te & 29.78 & 28.12 & 29.77 \\
linzhi valley    & Tr & 28.91 & 29.62 & 28.41 \\
bj south         & Va & 27.84 & 27.26 & 27.36 \\
hohhot steppe    & Tr & 26.15 & 27.18 & 26.55 \\
zhengzhou east   & Tr & 26.13 & 27.21 & 26.15 \\
dg north         & Te & 27.01 & 26.20 & 24.61 \\
sz east          & Te & 26.33 & 25.82 & 25.24 \\
haikou north     & Te & 26.48 & 25.72 & 24.95 \\
kunming plateau  & Va & 25.72 & 25.18 & 24.74 \\
urumqi north     & Tr & 24.68 & 25.77 & 24.84 \\
xiamen inland    & Te & 25.70 & 24.69 & 23.59 \\
wuhan lake       & Tr & 24.15 & 25.36 & 23.76 \\
harbin agri      & Tr & 23.51 & 24.19 & 23.49 \\
qingdao inland   & Va & 18.64 & 18.23 & 18.82 \\
\midrule
Range            &    & \textbf{12.45} & \textbf{13.08} & \textbf{10.96} \\
Train mean (10)  &    & 27.29 & 28.20 & 26.51 \\
Val mean (3)     &    & 24.07 & 23.56 & 23.64 \\
Test mean (5)    &    & 27.06 & 26.11 & 25.63 \\
\bottomrule
\end{tabular}
\end{table}

\subsection{Scene-Level View}

The same comparison at scene level, over the 1314 samples that retain a non-empty
valid mask for both networks, gives Mamba-MISR the better score in 36.6\% of scenes
and a 0.195~dB loss on average (Fig.~\ref{fig:scene}). The direction and magnitude
agree with the site-level result. Stratifying scenes by the PSNR of Mamba-MISR does
not explain the pattern, since the win rate by quartile, from hardest to easiest,
is 0.343, 0.366, 0.424 and 0.331 and the correlation with scene difficulty is 0.17.

\begin{figure}[!t]
\centering
\includegraphics[width=\columnwidth]{fig3_scene_winrate.pdf}
\caption{(a) Per-scene PSNR difference against the PSNR of Mamba-MISR, with
quartile means overlaid. (b) Mamba-MISR win rate by difficulty quartile. The
advantage is not concentrated in easy or hard scenes.}
\label{fig:scene}
\end{figure}

\subsection{How Many Sites Are Enough, and External Replication}
\label{sec:repl}

Fig.~\ref{fig:flip} also traces the win rate as the subset size grows. The values
for Mamba-MISR against HighRes-Net-cld are 0.390 at $k = 3$, 0.329 at $k = 5$,
0.261 at $k = 8$, 0.213 at $k = 10$, 0.154 at $k = 12$, 0.069 at $k = 15$ and
0.000 at $k = 16$. The win rate crosses the 5\% threshold only at $k = 16$. For a
gap of this size twelve sites would still leave a 15\% chance of reporting the
wrong winner, and the five-site test sets in routine use sit far outside the stable
region.

The replication uses a benchmark built by others, from a different sensor pair, at
a different scale factor, and with no model training at all. Under Setting~A the
PSNR range across the seven MGRS tiles is 5.54~dB for nearest-neighbour, 6.92~dB
for bicubic and 6.97~dB for bilinear, as Fig.~\ref{fig:mus2}(a) shows, and the
pattern holds at both wavelengths tested. Averaged over the six Setting~A baselines
the range is 6.31~dB. Under Setting~B the range is 14.13 and 8.77~dB for the two
bands, averaging 11.45~dB.

Fig.~\ref{fig:mus2}(b) places these ranges beside the S2DS value of 11.38~dB and
beside the 0.05 to 0.20~dB band occupied by typical method gaps. Setting~A, which
carries no registration uncertainty, sits at 32 times the upper end of that band,
and Setting~B and S2DS sit at 57 times.

Site variance on MuS2 has a standard deviation of 2.87~dB in band b3 and 2.70~dB in
b6, closely matching the 2.81~dB of S2DS. The two bands are analysed separately
because they differ by 5 to 8~dB in absolute level.

\begin{figure}[!t]
\centering
\includegraphics[width=\columnwidth]{fig4_mus2.pdf}
\caption{(a) Self-consistent PSNR of three zero-training baselines at each of the
seven MuS2 tiles. (b) Range across sites for S2DS, MuS2 Setting~A and MuS2
Setting~B, against the 0.05--0.20~dB band occupied by typical method gaps.}
\label{fig:mus2}
\end{figure}

\subsection{Valid-Pixel Fraction}

A second and independent way a benchmark can mislead emerged during the site-level
analysis. Fig.~\ref{fig:valid} shows that the fraction of valid, cloud-free pixels
in the target image varies by a factor of seven across the 18 sites, from 0.147 to
0.994. The scene-level distribution is more extreme. At \texttt{qingdao\_inland}
the median valid fraction is 0.015, so more than half of its scenes have almost no
usable pixels, and its PSNR is the lowest of the 18 sites at 18.56~dB. The
relationship between valid fraction and PSNR is positive but partial, with a
Pearson correlation across sites of 0.54, and \texttt{qingdao\_inland} sits six to
eight decibels below the trend. The valid fraction accounts for part of the site's
weakness and not all of it, and the remaining cause is not identified here. What
matters for reporting is that a mean over 1.5\% of a scene carries a variance that
no number of scenes removes when the same acquisition conditions recur.

\begin{figure}[!t]
\centering
\includegraphics[width=\columnwidth]{fig5_validfrac.pdf}
\caption{Valid (cloud-free) pixel fraction against valid PSNR for the 18 sites.
The relationship is positive but partial, and the lowest-scoring site sits well
below the trend.}
\label{fig:valid}
\end{figure}

%% ---------------------------------------------------------------- IV
\section{Discussion}

The two strongest networks in this study differ by 0.16~dB. Their site-level paired
difference is not distinguishable from zero, and over five-site test subsets the
ordering reverses in one case out of three. The same networks vary by 11.38~dB
across sites, and the mean score moves by 6.54~dB depending on which five sites are
drawn. The last two figures are one to two orders of magnitude larger than the
quantity being resolved. The fragile quantity is the narrow gap, which is the
regime that method papers occupy when they report an improvement.

A benchmark score has always been understood to depend on its test data. The
quantities reported here give that dependence a magnitude, specific to a benchmark
family and to a gap size. How much does the score move when the test sites change.
How many sites are needed before a ranking of a given gap size settles. At what gap
size does the ranking stop depending on site composition. The answers here are
6.54~dB, 16 sites, and somewhere between 0.16 and 0.90~dB. The MuS2 replication
matters for the same reason, since the effect appears on a benchmark built by other
authors with different sensors and a different scale factor, and the magnitude
differs in the direction that the two benchmarks' geographical extents predict.

Four reporting practices follow. Report per-site scores, since a single aggregate
conceals a distribution whose spread exceeds the effect being claimed, and the
per-site values cost a few lines of a table. Run leave-one-site-out, which shows
whether the headline number survives the removal of any single site. Report rank
stability rather than a single ordering, since for two methods within a decibel of
each other a win rate over plausible test sets carries more information than a
rank. Treat the number of sites as a design parameter, because a test set of three
to five geographically distinct sites cannot resolve differences below roughly one
decibel on this benchmark. None of the four requires extra training.

%% ---------------------------------------------------------------- V
\section{Conclusion}

A benchmark score in remote sensing SR is decided to a large extent by which sites
the benchmark happens to contain. On a multi-temporal Sentinel-2 benchmark, methods
differing by 0.16~dB exchange places in one third of five-site test sets,
site-to-site variance exceeds model-to-model variance by 31.7 times, and the
benchmark's own three partitions yield three different orderings of the same
models. At least 16 sites are needed before a gap of this size becomes resolvable.
The effect reproduces on an independently built public benchmark where no model is
trained at all. Reporting per-site scores, running leave-one-site-out, and stating
the rank stability of narrow gaps would make these limits visible at no
experimental cost.

%% References are ordered by first citation, as IEEE requires.
\begin{thebibliography}{12}

\bibitem{deepsum2020}
A.~B. Molini, D.~Valsesia, G.~Fracastoro, and E.~Magli, ``DeepSUM: Deep neural
network for super-resolution of unregistered multitemporal images,'' \emph{IEEE
Trans. Geosci. Remote Sens.}, vol.~58, no.~5, pp.~3644--3656, May 2020.

\bibitem{diffusionsat2024}
S.~Khanna, P.~Liu, L.~Zhou, C.~Meng, R.~Rombach, M.~Burke, D.~B. Lobell, and
S.~Ermon, ``DiffusionSat: A generative foundation model for satellite imagery,''
in \emph{Proc. Int. Conf. Learning Representations (ICLR)}, 2024.

\bibitem{mus2}
P.~Kowaleczko \emph{et al.}, ``MuS2: A real-world benchmark for Sentinel-2
multi-image super-resolution,'' \emph{Sci. Data}, vol.~10, 2023, Art.~712.

\bibitem{drusch2012}
M.~Drusch, U.~Del Bello, S.~Carlier, O.~Colin, V.~Fernandez, F.~Gascon,
\emph{et al.}, ``Sentinel-2: ESA's optical high-resolution mission for GMES
operational services,'' \emph{Remote Sens. Environ.}, vol.~120, pp.~25--36,
May 2012.

\bibitem{s2psd}
ESA, \emph{Sentinel-2 Products Specification Document}, issue 14.9, 2023.

\bibitem{mamba}
A.~Gu and T.~Dao, ``Mamba: Linear-time sequence modeling with selective state
spaces,'' arXiv:2312.00752, Dec.~2023.

\bibitem{deudon2020}
M.~Deudon, A.~Kalaitzis, S.~Gowda, \emph{et al.}, ``HighRes-net: Recursive fusion
for multi-frame super-resolution of satellite imagery,'' in \emph{Proc. IEEE/CVF
CVPR Workshops}, 2020, pp.~298--307.

\bibitem{okabayashi2024}
A.~Okabayashi, N.~Audebert, S.~Donike, and C.~Pelletier, ``Cross-sensor
super-resolution of irregularly sampled Sentinel-2 time series,'' in \emph{Proc.
IEEE/CVF CVPR Workshops (EarthVision)}, 2024, pp.~502--511.

\bibitem{cmix}
S.~Skakun, J.~Wevers, C.~Brockmann, G.~Doxani, \emph{et al.}, ``Cloud Mask
Intercomparison eXercise (CMIX): An evaluation of cloud masking algorithms for
Landsat~8 and Sentinel-2,'' \emph{Remote Sens. Environ.}, vol.~274, 2022,
Art.~112990.

\bibitem{ssim2004}
Z.~Wang, A.~C. Bovik, H.~R. Sheikh, and E.~P. Simoncelli, ``Image quality
assessment: from error visibility to structural similarity,'' \emph{IEEE Trans.
Image Process.}, vol.~13, no.~4, pp.~600--612, Apr.~2004.

\bibitem{wilcoxon1945}
F.~Wilcoxon, ``Individual comparisons by ranking methods,'' \emph{Biometrics
Bulletin}, vol.~1, no.~6, pp.~80--83, Dec.~1945.

\bibitem{lakens2017}
D.~Lakens, ``Equivalence tests: A practical primer for $t$ tests, correlations,
and meta-analyses,'' \emph{Soc. Psychol. Personal. Sci.}, vol.~8, no.~4,
pp.~355--362, 2017.

\end{thebibliography}

%% Title-page footnote material moved here (end of the manuscript, last page).
\vfill
\noindent{\footnotesize\raggedright Manuscript received XXXX; revised XXXX. The
authors are with the Engineering University of PAP, Xi'an~710000, China.\\
$^*$Corresponding author: Dawei Liu (e-mail: wjmicheal@163.com).\\
Code and the per-site evaluation tables are available at
\url{https://github.com/[[FILL-REPO]]}.}

\end{document}
